

\section{Fermat's Little Theorem}
Fermat's Little Theorem is one of the oldest results in elementary number theory, recorded in a 1640 letter from Pierre de Fermat to Fr\'{e}nicle de Bessy. It underlies essentially every computation we perform in the multiplicative group $\bbF_p^*$ -- from the very existence of the canonical additive character of $\bbF_q$ (where $q=p^r$), through the character congruence $\psi(l)\equiv l^f\pmod P$ used in the arithmetic proof of the Davenport--Hasse relation, down to the mod-$q$ reduction of Gauss sums appearing in the proof of the \QRthm. We record it here for completeness.
\begin{Theorem}{Fermat's Little Theorem}{thm:fermat-little}
  Let $p$ be a prime. Then for every integer $a$ we have $$a^p\equiv a\pmod p,$$ and in particular, if $p\nmid a$ then $a^{p-1}\equiv 1\pmod p$.
\end{Theorem}
\begin{proof}
  The non-zero residues modulo $p$ form a group $\bbF_p^*=\{1,2,\dots,p-1\}$ under multiplication of order $p-1$. By Lagrange's theorem every element of a finite group satisfies $g^{|G|}=e$, so $a^{p-1}\equiv 1\pmod p$ for every $a\in\bbF_p^*$, i.e.\ every $a$ with $p\nmid a$. Multiplying both sides by $a$ gives $a^p\equiv a\pmod p$ for such $a$; and the congruence $a^p\equiv a\pmod p$ is trivially true when $p\mid a$, so it holds for every integer $a$.
\end{proof}

\section{Waring's Formula}
Waring's Formula expresses each power-sum symmetric polynomial $s_k=x_1^k+x_2^k+\cdots+x_n^k$ as an explicit polynomial in the elementary symmetric polynomials $\sigma_1,\dots,\sigma_n$. It is the source of the closed-form reduction of lifted Kloosterman sums $K(\chi^{(s)};a,b)$ to the ground-field sum $K(\chi;a,b)$ in \thmref{thm:kloosterman-reduction}: applying Waring's formula to $\omega_1^s+\omega_2^s$ with $\sigma_1=\omega_1+\omega_2=-K(\chi;a,b)$ and $\sigma_2=\omega_1\omega_2=q$ produces the binomial-coefficient expansion. We state the general identity here for reference.
\begin{Theorem}{Waring's Formula}{thm:waring}
  Let $x_1,\dots,x_n$ be indeterminates, let $$\sigma_j=\sum_{1\leq i_1<i_2<\cdots<i_j\leq n}x_{i_1}x_{i_2}\cdots x_{i_j}\qquad (j=1,\dots,n)$$ be the elementary symmetric polynomials in $x_1,\dots,x_n$, and let $$s_k=x_1^k+x_2^k+\cdots+x_n^k$$ be the $k^{\text{th}}$ power-sum. Then for every integer $k\geq 1$,
  $$s_k=\sum(-1)^{i_2+i_4+i_6+\cdots}\cdot \frac{(i_1+i_2+\cdots+i_n-1)!\cdot k}{i_1!\cdot i_2!\cdots i_n!}\cdot \sigma_1^{i_1}\sigma_2^{i_2}\cdots \sigma_n^{i_n},$$
  where the summation runs over all $n$-tuples $(i_1,\dots,i_n)$ of non-negative integers satisfying $i_1+2i_2+\cdots+ni_n=k$. Each coefficient of $\sigma_1^{i_1}\sigma_2^{i_2}\cdots \sigma_n^{i_n}$ is an integer.
\end{Theorem}